\documentclass[../../../include/open-logic-section]{subfiles}

\begin{document}

\olfileid{cmp}{thy}{clo}

\olsection[可计算枚举集的并与交]{可计算枚举集在并与交运算下封闭}
% Part: computability
% Chapter: computability-theory
% Section: ce-closed-cup-cap

下面的定理给出了可计算枚举集合的一些封闭性质。

\begin{thm}
设 $A$ 和 $B$ 是可计算枚举的。那么 $A \cap B$ 和 $A \cup B$ 也是可计算枚举的。
\end{thm}

\begin{proof}
根据 \olref[eqc]{thm:ce-equiv}，我们可以使用可计算枚举集合的各种等价刻画。作为示例，我们将提供几种不同的证明。

第一种证明：假设 $A$ 由可计算函数~$f$ 枚举，$B$ 由可计算函数~$g$ 枚举。令
\begin{align*}
h(x) & = \umin{y}{(f(y) = x \lor g(y) = x)} \text{且}\\
j(x) & = \umin{y}{(f((y)_0) = x \land g((y)_1) = x)}.
\end{align*}
则 $A \cup B$ 是 $h$ 的定义域，$A \cap B$ 是~$j$ 的定义域。

\begin{explain}
计算上的直观是这样的：给定枚举 $A$ 和 $B$ 的过程，我们可以通过在两个枚举中的任意一个里寻找 $x$，来半判定元素 $x$ 是否在 $A \cup B$ 中；而通过同时在两个枚举中寻找 $x$，我们可以半判定元素 $x$ 是否在 $A \cap B$ 中。
\end{explain}

第二种证明：再次假设 $A$ 由~$f$ 枚举，$B$ 由~$g$ 枚举。令
\[
k(x) = \begin{cases}
f(x/2) & \text{若 $x$ 是偶数} \\
g((x-1)/2) & \text{若 $x$ 是奇数。}
\end{cases}
\]
则 $k$ 枚举了 $A \cup B$；其思路是 $k$ 只是在 $f$ 和~$g$ 提供的枚举之间交替进行。
枚举 $A \cap B$ 则更棘手一些。如果 $A \cap B$ 是空集，那它显然是可计算枚举的。
否则，令 $c$ 为 $A \cap B$ 的任意元素，并定义 $l$ 如下：
\[
l(x) = \begin{cases}
f((x)_0) & \text{若 $f((x)_0) = g((x)_1)$} \\
c & \text{否则。}
\end{cases}
\]
从计算的角度看，$l$ 遍历 $f$ 和 $g$ 的枚举中的元素对，并输出找到的每一个匹配；若未找到匹配，它只是通过输出 $c$ 来“原地踏步”。

最后一种证明：假设 $A$ 是偏函数 $m(x)$ 的\emph{定义域}，$B$ 是偏函数 $n(x)$ 的定义域。
那么 $A \cap B$ 就是偏函数 $m(x) + n(x)$ 的定义域。

\begin{explain}
用计算术语来说，如果 $A$ 是使 $m$ 停机的值的集合，$B$ 是使 $n$ 停机的值的集合，
那么 $A \cap B$ 就是使两个过程都停机的值的集合。
\end{explain}

将 $A \cup B$ 表示为停机值的集合则更为困难，因为必须并行模拟 $m$ 和 $n$。
令 $d$ 是 $m$ 的一个指标，$e$ 是 $n$ 的一个指标；换言之，$m = \cfind{d}$ 且 $n = \cfind{e}$。
那么 $A \cup B$ 就是如下函数的定义域：
\[
p(x) = \umin{y}{(T(d,x,y) \lor T(e,x,y))}.
\]
\begin{explain}
从计算的角度看，在输入 $x$ 下，$p$ 会搜索 $m$ 停机的计算或者 $n$ 停机的计算，
并在找到任意一个时停机。
\end{explain}
\end{proof}
\end{document}